\scanpage{10}
\mspage{8}

\uline{à isogénie près}'', ou \uline{isogristaux}, sur S. Utilisant cette nouvelle notion, on définit comme ci-dessus la notion de p-isocristal, comme étant la donnée d'un isocristal $\underline M$ muni d'un homomorphisme $F:\underline M^{(p)}\longrightarrow\underline M$. La notion de bi-isocristal de poids i, qui serait calquée de celle de bi-cristal de poids i, n'est pas très raisonnable alors, car V doit être alors donné en termes de F comme $p^{\deleted{\unsure{2}}i}F^{-1}$. Il y a lieu plutôt d'appeler \uline{bi-isocristal} (sans précision de poids) un p-isocristal pour lequel F est un \uline{isomorphisme}, et de ne pas choisir entre les différents $p^{\deleted{\unsure{2}}i}F^{-1}$ possibles; il vaut mieux également ne pas parler ici de $p^{-1}$-isocristal, la notion intéressante étant celle de bi-isocristal, qui est manifestement \uline{autoduale} \deleted{(si $\unsure{\phi|F}$ est un foncteur contravariant de catégories fibrées sur un} \inserted{quand on se restreint aux objets de type fini:} le dual de $(\underline M,F)$ est $(\underline M,{}^tF^{-1})$, où $\underline M$ est le iso-cristal dual de M.

On peut également, bien sûr, localiser par rapport à p en partant déjà de la catégorie des bi-cristaux de poids i, on trouve les \uline{bicristaux de poids} i \uline{à isogénie près}, qui forment une catégorie abélienne $\mathrm{Is}\inserted{\mathrm{bi}}\mathrm{cr}(S,i)$, et un foncteur exact ``oubli de V''
\[
\unsure{\operatorname{Isbicr}|\operatorname{IsBicr}}(\unsure{S,i|S)i})\longrightarrow\operatorname{Isbicr}(S),
\]
à valeurs dans la catégorie $\operatorname{Isbicr}(S)$ des iso-bi-cristaux sur S. Ce foncteur est \deleted{\unsure{évidemment}} \uline{pleinement fidèle}, de sorte que \uline{les \inserted{\unsure{iso-}}bicristaux forment une généralisation commune des bicristaux de poids} i \uline{à isogénie près}. Pour préciser ce point, désignant par $\operatorname{Bicr}(S,i)$ la catégorie des bicristaux de poids i sur S, il y a lieu d'introduire des foncteurs canoniques
\[
\operatorname{Bicr}(S,i)\longrightarrow\operatorname{Bicr}(S,i+1)\longrightarrow\ldots,
\]
donnés par $(\underline M,F,V)\rightsquigarrow(\underline M,F,pV)$. Quand on localise ces foncteurs par p, on trouve des foncteurs
\[
\operatorname{IsBicr}(S,i)\longrightarrow\operatorname{IsBicr}(S,i+1)\longrightarrow\ldots.
\]
qui se trouvent être \uline{pleinement fidèles} : on peut donc considérer que, pour i croissant, la notion de ``bicristal de poids i, à isogénie près'' constitue une généralisation de plus en plus large de celle de ``cristal de Dieudonné''. La noti-
\margin{9}
